%==========================================================================
% APPENDICES
%==========================================================================
\appendix

\chapter{Code Repository}
\label{app:code}

The complete source code for this project is available at the following repository:

\begin{center}
\url{https://github.com/Rivi9/dynamic-quest-gen}
\end{center}

The repository contains:
\begin{itemize}
\item \texttt{backend/} --- FastAPI Python backend (player modelling, narrative generation, RL adaptation)
\item \texttt{unity-plugin/} --- UPM package (telemetry collection, narrative management, content injection)
\item \texttt{demo-game/} --- Unity 6 testbed project
\item \texttt{lore/} --- Game world lore documents (RAG corpus)
\item \texttt{evaluation/} --- Questionnaires, interview guide, and statistical analysis scripts
\item \texttt{data/} --- Trained PPO model weights and SQLite database
\end{itemize}

\textit{Note: The full source code is not reproduced in this thesis per submission guidelines. Please refer to the repository for implementation details.}


\chapter{GEQ Core Module}
\label{app:geq}

The Game Experience Questionnaire Core Module as administered in this study. Reproduced from \citet{ijsselsteijn2013geq} for research purposes.

\textbf{Scale:} 0 = Not at all, 1 = Slightly, 2 = Moderately, 3 = Fairly, 4 = Very much

\textit{Please answer each question based on your experience during the game session you just played.}

\begin{longtable}{@{}clc@{}}
\toprule
\textbf{\#} & \textbf{Item} & \textbf{Scale} \\
\midrule
\endfirsthead
\toprule
\textbf{\#} & \textbf{Item} & \textbf{Scale} \\
\midrule
\endhead
1 & I felt content & 0--4 \\
2 & I felt skillful & 0--4 \\
3 & I was absorbed in the game & 0--4 \\
4 & I felt happy & 0--4 \\
5 & This game gave me a bad mood & 0--4 (R) \\
6 & I thought it was fun & 0--4 \\
7 & I was fully occupied with the game & 0--4 \\
8 & I felt annoyed & 0--4 (R) \\
9 & I found the game satisfying & 0--4 \\
10 & I felt frustrated & 0--4 (R) \\
11 & It felt like a whole new world & 0--4 \\
12 & I lost track of time & 0--4 \\
13 & I felt bored & 0--4 (R) \\
14 & I felt good & 0--4 \\
15 & I was good at this game & 0--4 \\
16 & I felt tired & 0--4 (R) \\
17 & I felt imaginative & 0--4 \\
18 & I felt that I could explore things & 0--4 \\
19 & I felt like I had to do things quickly & 0--4 \\
20 & I had fun with others & 0--4 \\
21 & I felt challenged & 0--4 \\
22 & I felt that time went by quickly & 0--4 \\
23 & I felt pressured & 0--4 (R) \\
24 & I felt competent & 0--4 \\
25 & I felt engrossed & 0--4 \\
26 & I felt bothered by aspects of the game & 0--4 (R) \\
27 & I was deeply concentrated in the game & 0--4 \\
28 & I forgot everything around me & 0--4 \\
29 & I felt pleased & 0--4 \\
30 & I felt completely focused on the game & 0--4 \\
31 & I enjoyed the game & 0--4 \\
32 & I felt imaginative & 0--4 \\
33 & I felt adventurous & 0--4 \\
\bottomrule
\end{longtable}

\textit{(R) = reverse-scored}

\vspace{0.5cm}
\textbf{Subscale scoring:}

\begin{table}[H]
\centering
\small
\begin{tabular}{@{}lll@{}}
\toprule
\textbf{Subscale} & \textbf{Items} & \textbf{Notes} \\
\midrule
Flow (Primary DV, H1) & 5(R), 13(R), 25, 28, 31 & Mean of 5 items \\
Immersion (Secondary DV, H2) & 3, 12, 18, 19, 27, 30 & Mean of 6 items \\
Positive Affect & 1, 4, 8(R), 20, 29 & Mean of 5 items \\
Negative Affect & 10, 16, 23, 26 & Mean of 4 items \\
Tension & 7, 22, 25 & Mean of 3 items \\
Competence & 2, 15, 17, 21, 24 & Mean of 5 items \\
\bottomrule
\end{tabular}
\end{table}


\chapter{miniPXI Questionnaire}
\label{app:minipxi}

Citation: \citet{vornhagen2022minipxi}

\textbf{Scale:} 1 = Strongly Disagree \ldots{} 7 = Strongly Agree

\begin{longtable}{@{}cl@{}}
\toprule
\textbf{\#} & \textbf{Item} \\
\midrule
\endfirsthead
1 & I felt in control of my character/avatar \\
2 & I felt a sense of mastery when completing challenges \\
3 & I found the game world interesting to explore \\
4 & I was curious about what would happen next in the game \\
5 & I felt like the game narrative pulled me in \\
6 & I had a sense of presence in the game world \\
7 & I felt free to make meaningful choices \\
8 & The game felt meaningful to me \\
9 & I had positive emotions while playing \\
10 & I felt excited during the game \\
11 & Overall, I enjoyed playing this game \\
\bottomrule
\end{longtable}

\textbf{Subscales:} Mastery (1, 2); Curiosity (3, 4); Immersion (5, 6); Autonomy (7, 8); Positive Affect (9, 10, 11). Total score: mean of all 11 items.


\chapter{Narrative Personalization Scale (NPS-3)}
\label{app:nps3}

The NPS-3 is a custom three-item scale developed for this study to measure perceived narrative personalization. Cronbach's alpha will be reported in Chapter~5.

\textbf{Scale:} 1 = Strongly Disagree \ldots{} 7 = Strongly Agree

\begin{longtable}{@{}cl@{}}
\toprule
\textbf{\#} & \textbf{Item} \\
\midrule
\endfirsthead
1 & The story felt tailored to how I was playing. \\
2 & The narrative responded to my choices and actions. \\
3 & The game seemed to understand what kind of player I am. \\
\bottomrule
\end{longtable}

\textbf{Scoring:} Mean of all 3 items. If Cronbach's $\alpha < 0.70$, items are analysed individually.


\chapter{Interview Guide}
\label{app:interview}

\textbf{Study:} AI-Driven Dynamic Narrative Personalization in Games \\
\textbf{Duration:} $\sim$20 minutes \\
\textbf{Analysis method:} Reflexive Thematic Analysis \citep{braun2022thematic}

\textbf{Before the interview:} Obtain signed consent; remind participant there are no right or wrong answers; start audio recording.

\vspace{0.5cm}
\textbf{Section 1 --- General Experience (5 min)}
\begin{enumerate}
\item Can you walk me through what you remember most about the game experience?
\item How would you describe the story and the way it unfolded for you?
\item Were there any moments that stood out to you, positively or negatively?
\end{enumerate}

\textbf{Section 2 --- Narrative and Personalization (8 min)}
\begin{enumerate}[resume]
\item Did the narrative feel connected to the way you were playing? Can you give an example?
\item Were there moments where the story surprised you --- or felt unexpectedly relevant?
\item Did any NPCs or dialogue feel particularly fitting or particularly out of place?
\item \textit{(Experimental only)} Did you notice the story adapting to you? How did that feel?
\item \textit{(Control only)} How well did the story match your play style?
\end{enumerate}

\textbf{Section 3 --- Flow and Engagement (5 min)}
\begin{enumerate}[resume]
\item Were there moments where you felt in ``the zone'' --- neither too stressed nor too bored?
\item Were there moments where the game felt too easy, too hard, or just right?
\item How did the story contribute to those moments of engagement?
\end{enumerate}

\textbf{Section 4 --- Closing (2 min)}
\begin{enumerate}[resume]
\item Is there anything about the game or narrative you would like to add?
\item Any suggestions for improving the experience?
\end{enumerate}

\textbf{Probes:} ``Can you say more about that?'' / ``What do you mean when you say [X]?'' / ``When did that happen in the game?'' / ``How did that make you feel?''


\chapter{API Specification}
\label{app:api}

The FastAPI backend exposes the following REST and WebSocket endpoints:

\begin{table}[H]
\centering
\small
\begin{tabular}{@{}llp{6cm}@{}}
\toprule
\textbf{Method} & \textbf{Path} & \textbf{Description} \\
\midrule
POST & /api/telemetry & Submit telemetry window; persists batch to SQLite \\
POST & /api/narrative/generate & Generate narrative content for a session \\
GET & /api/player-model/\{session\_id\} & Get current PlayerModel (Hexad + Flow) \\
GET & /health & Health check \\
WebSocket & /ws/telemetry & Low-latency telemetry stream \\
\bottomrule
\end{tabular}
\end{table}

\textbf{NarrativeContent response example:}

\begin{lstlisting}[language={}]
{
  "action_taken": "PROVIDE_GUIDANCE",
  "content_type": "dialogue",
  "content": "The path below is treacherous. If you
    still carry the torch, keep it lit.",
  "speaker": "Commander Varis",
  "emotional_tone": "tense",
  "lore_refs": [],
  "fallback": false
}
\end{lstlisting}

All endpoints return standard HTTP status codes. 422 for malformed requests (Pydantic validation). 503 if Ollama is unreachable (returns fallback NarrativeContent with \texttt{"fallback": true}).


\chapter{MDP Formal Specification}
\label{app:mdp}

\textbf{Formal MDP tuple:} $(S, A, T, R, \gamma)$

\textbf{State space $S$:} Box(7,) $\in [0, 1]^7$

\begin{table}[H]
\centering
\small
\begin{tabular}{@{}clcl@{}}
\toprule
\textbf{Dim} & \textbf{Symbol} & \textbf{Range} & \textbf{Description} \\
\midrule
$s_0$ & flow\_score & \{0.1, 0.2, 0.3, 1.0\} & Encoded Flow state \\
$s_1$ & csr\_norm & [0, 1] & challenge\_skill\_ratio / 3.0 \\
$s_2$ & $h_{\text{explorer}}$ & [0, 1] & Hexad Explorer score \\
$s_3$ & $h_{\text{socializer}}$ & [0, 1] & Hexad Socializer score \\
$s_4$ & $h_{\text{achiever}}$ & [0, 1] & Hexad Achiever score \\
$s_5$ & $h_{\text{disruptor}}$ & [0, 1] & Hexad Disruptor score \\
$s_6$ & $t_{\text{progress}}$ & [0, 1] & session\_elapsed / 1800 \\
\bottomrule
\end{tabular}
\end{table}

\textbf{Reward function $R(s, a, s')$:}

\[
R = R_{\text{state}}(s') + R_{\text{step}} + R_{\text{streak}}(s', n) + R_{\text{transition}}(s, s') + R_{\text{repetition}}(a, a_{[-3:]})
\]

Where:
\begin{itemize}
\item $R_{\text{state}}$: FLOW $= +1.0$, BOREDOM $= -0.3$, ANXIETY $= -0.5$, APATHY $= -0.8$
\item $R_{\text{step}} = -0.05$
\item $R_{\text{streak}} = +0.20$ if $s' = \text{FLOW}$ and $n > 2$
\item $R_{\text{transition}} = +0.30$ if $s \in \{\text{ANXIETY}, \text{BOREDOM}\}$ and $s' = \text{FLOW}$
\item $R_{\text{repetition}} = -0.15$ if all last 3 actions identical
\end{itemize}

\textbf{Discount factor $\gamma$:} 0.95

\textbf{Episode length:} 360 steps (30 minutes at 5-second intervals)

\textbf{Training:} PPO (Stable-Baselines3 v2.x), 200,000 timesteps, 4 parallel environments.
